Here we continue on following Witten's discussion on the relation
between 2+1 dimensional gravity and Chern-Simons theory
\cite{witten21Dimensional1988}.
\subsection{Geometric Phase Space Analysis}

\YL{[Plan to not keep this , its unclear!!]}

\subsubsection{Discussion Set Up}

How to quantize a Theory, we usually start from analyzing the phase
space of the theory and then construct a poisson braket of the phase
space. Then one can transform the poisson bracket into commutator and
get the quantum theory.
Thus to study the quantum gravity in 2+1 dimensions, we need to first
learn about the phase space of the classical 2+1 dimensional gravity.

The phase space of a classical theory is often defined as:
\defi{
  Phase Space

  The phase space is the space of all $ q, \dot{q} $ at $ t = 0 $ surface.
}
However, this definition is not covariant and we need tons of
constraints to make it work, thus in 2+1 dimensional gravity, we use
the following definition:
\defi{
  Phase Space

  The phase space is the space of all classical solutions to the
  equations of motion Modulo gauge transformations.
}
For 2+1 dimensional gravity, this means that we want to find all the
solutions to:
\begin{align}
  R_{\mu\nu} = 0
\end{align}
To do this, we might first assume a spacetime topology of the
manifold we want to consider, we focus on a manifold:
\begin{itemize}
  \item With topology $ M \cong [0,1] \times \Sigma $. where $ \Sigma
    $ is a genus $ g $ closed surface.
\end{itemize}

We now simplify the problem of analyzing the phase space of 2+1
dimensional gravity to finding the space of all solutions to the
vacuum Einstein equation $ R_{\mu\nu} = 0 $ on a manifold with
topology $ M \cong [0,1] \times \Sigma $. These solutions can be
found by quotienting the Minkowski space by a discrete subgroup of
the isometry group of Minkowski space, which is the Poincare group $
ISO(2,1) $. In Witten's paper he gives out a analysis of those manifolds.

\subsubsection{Manifolds with initial \& final singularities}

Witten explicitly constructs these class of solutions as examples, by
quotienting
the light cone of Minkowski, now we explain the procedure:

\begin{itemize}
  \item \textbf{Preparation: Quotienting $ \mathbb{H} $ to $ \Sigma_g $}

    We can construct a 2-manifold with genus $ g \geq 2 $ by
    quotienting the hyperbolic plane.
    \begin{itemize}
      \item \textbf{Step 1: } We first find the fundamental group of
        the surface $ \Sigma_g $ which is $ \pi_1(\Sigma_g) $.
      \item \textbf{Step 2: } We find a subgroup of $
        PSL(2,\mathbb{R}) $ naming $ \Gamma $
        which is isomorphic to $ \pi_1(\Sigma_g) $ and acts
        \textbf{discretely} on $ \mathbb{H} $.
      \item \textbf{Step 3: } The quotient $ \mathbb{H}/\Gamma  $
        defines a Riemann surface with genus $ g $ and a metric with
        constant negative curvature.
    \end{itemize}
\end{itemize}
\YL{[In fact I don't understand how this step works mathematically, eg
    how to construct the metric and why we don't require the group to
    be free but only properly continuous (discrete), etc. I need to learn
more about this step.]}
\rmk{
  Note that here we quotient by a discrete subgroup other than a
  properly discontinuous subgroup, because the action of a properly
  discontinuous subgroup is equivalent to being discrete in the case
  of being a subgroup of $ Aut(\mathbb{H}) = PSL(2,\mathbb{R}) $.
}

Then we apply a similar procedure to the light cone of Minkowski space to
construct a flat 3-manifold with topology $ M \cong [0,1] \times
\Sigma_g $ with initial singularity and final singularity.

To do this we first list out some useful facts about the 2+1
Minkowski spacetime:
\begin{itemize}
  \item Normally we have the metric of:
    \begin{align}
      ds^2 = -dt^2 + dx^2 + dy^2
    \end{align}
    And the light cone is defined by:
    \begin{itemize}
      \item $ X^+ $ future light cone: $ t^2 > x^2 + y^2, t > 0 $
      \item $ X^- $ past light cone: $ t^2 > x^2 + y^2, t < 0 $
    \end{itemize}
  \item The isometry group is the 3 dimensional Poincare group $
    ISO(2,1) $ which is the semidirect product of
    the Lorentz group $ SO(2,1) $ and the translation group $
    \mathbb{R}^{2,1} $.
    \begin{itemize}
      \item The proper orthochronous Lorentz group $ SO^+(2,1) $ is
        isomorphic to $ PSL(2,\mathbb{R}) $.
    \end{itemize}
\end{itemize}
\textbf{Quotienting the light cone to get a 3-manifold with initial \&
final singularities}

We first notice that the light cone can be made up by hyperbolic planes:
\begin{align}
  t^2 = x^2 + y^2 + \tau^2
\end{align}
where $ \tau $ is a parameter for the hyperbolic plane.
\begin{itemize}
  \item Each hyperbolic plane is isometric to $ \mathbb{H} $ and the
    isometry group of the hyperbolic plane is the Lorentz group $
    SO^+(2,1) $ which is a subgroup of the Poincare group $ ISO(2,1) $
\end{itemize}
Thus given a
$ \Sigma_g $ we can construct a subgroup of the lorentz group $ \Gamma $
isometric to $ \pi_1(\Sigma_g) $ and quotient the hyperbolic plane by
this subgroup to get a $ \Sigma_g $. If we do this quotient for the
whole $ X^+ $ light cone we can get a 3 manifold with topology $ M
\cong [0,1] \times \Sigma_g $ and with an initial singularity at $ \tau
= 0 $.
\begin{itemize}
  \item The resulting manifold is a solution to the vaccum Einstein
    Field Equation for it is contructed by quotienting an isometry of
    Minkowski spacetime. This quotient group preserves the metric
    thus the resulting manifold is still flat.

    \YL{[I don't understand why?? It needs knowledge of metric under
    quotienting.]}
  \item The resulting manifold has a $ \Sigma_g \times [0,1] $
    topology. This is because the lorentz group act the hyperbolic
    plane to itself due to its feature of preserving Minkowski
    metric. Thus, $ \Gamma $ as a subgroup of lorentz group act
    independently on each hyperbolic
    plane.

    Moreover, the isometry group of the hyperbolic plane is
    the lorentz group which is isomorphic to $ PSL(2,\mathbb{R}) $
    and the hyperbolic plane is isometric to $
    \mathbb{H} $. Thus, quotienting the hyperbolic plane with the
    subgroup $ \Gamma $ of the lorentz group is equivalent to
    quotienting $ \mathbb{H} $ with the same subgroup, which gives us
    a $ \Sigma_g $. Thus, the resulting manifold has a $ \Sigma_g
    \times [0,1] $ topology.
\end{itemize}

\textbf{Parameters Counting}

Finally, we can count how many distinct manifold can be constructed
by this procedure, the answer depends on the choise of $ \Gamma $
subgroup. The final answer is that the moduli of genus $ g $ surface,
which is $ \mathbb{R}^{6g-6} $ gives us the label of distinct
manifolds we can construct by this procedure.

\YL{[I don't understand this counting, it need some knowledge of
Teichmuller space and moduli space of Riemann surface.]}

\subsubsection{General Flat Manifolds}

Now we try to construct more general flat manifolds with topology $ M
\cong [0,1] \times \Sigma_g $ by quotienting more general subspace of
the Minkowski space.

For a general flat 3-manifold $ \hat{M} $, we can construct its
simple connected universal cover $ M $. Due to the flatness and being
a universal covering space, $ M $ is isometric to the Minkowski space
or a subspace of the Minkowski space. And by definition, $ \hat{M} $
must be a quotient of $ M $, which is a subspace of the Minkowski
space, by a subgroup of the isometry group of the Minkowski space,
which is the Poincare group $ ISO(2,1) $.

Under this set up, we consider the fundamental group of $ \hat{M} $,
which is $ \pi_1(\hat{M}) $. Every element shall be a loop in $
\hat{M} $, which can be lifted to a path in $ M $. Then we can define
a element of the Poincare group $ ISO(2,1) $ $ \phi(\gamma) $ which
ensures that the
path $ [\gamma] $ in $ M $ which lifted from the loop $ [\gamma] $ in $
\hat{M} $, is mapped to a noncontractible loop after the quotienting
by $ \phi(\gamma) $.

Thus we can see that:
\begin{itemize}
  \item Every flat 3-manifold give rise to a homomorphism from $
    \pi_1(\hat{M}) $ to the Poincare group $
    ISO(2,1) $. It can be constructed by quotienting Minkowski space
    by the image of the homomorphism.
\end{itemize}

\YL{[but why this procedure can also cover quotienting a subspace of
    Minkowski space??? It only covers the case of quotienting the whole
    Minkowski space, but not the case of quotienting a subspace of
    Minkowski space, which is also a solution to the Einstein equation.
    I think perhaps this is one of the non-rigorous step in this
    derivation. Thus, we need a more rigorous field theoretical proof of
the phase space.]}

Thus, the phase space of 2+1 dimensional gravity can be identified
with the following space:
\begin{align}
  \text{Phase Space} = \text{Hom}(\pi_1(\Sigma_g), ISO(2,1))/\sim
\end{align}
Here $ \sim $ means that we quotient by the conjugation of the
Poincare group $ ISO(2,1) $. We know that if two homomorphisms $
\phi_1, \phi_2 $ are related as:
\begin{align}
  \phi_1(\gamma) = g \phi_2(\gamma) g^{-1} \quad \forall \gamma \in ISO(2,1)
\end{align}
Then they give rise to the same manifold under quotient.

If we want to consider the topology of $ \hat{M} = \Sigma_g \times [0,1] $,
then the fundamental group of the manifold is $ \pi_1(\hat{M}) =
\pi_1(\Sigma_g) $. Thus, the phase space of 2+1 dimensional gravity
with topology $ M \cong [0,1] \times \Sigma_g $ is:
\thm{
  The phase space of 2+1 dimensional gravity with topology $ M \cong
  [0,1] \times \Sigma_g $ is:
  \begin{align}
    \text{Phase Space} = \text{Hom}(\pi_1(\Sigma_g), ISO(2,1))/\sim
  \end{align}
  where $ \sim $ means that we quotient by the conjugation of the
  Poincare group $ ISO(2,1) $.
}

\textbf{Parameter Counting}

We then count the dimension of the phase space. For a surface $
\Sigma_g $, its fundamental group can be generated by $ 2g $
generators $ a_i, b_i $ with one relation:
\begin{align}
  \pi_1(\Sigma_g) = \langle a_1, b_1, \cdots, a_g, b_g |
  \prod_{i=1}^g a_i b_i a_i^{-1} b_i^{-1} = 1 \rangle
\end{align}
Thus a homomorphism from $ \pi_1(\Sigma_g) $ to the Poincare group
$ ISO(2,1) $ is determined by the image of the generators $ a_i, b_i $
in the Poincare group $ ISO(2,1) $. Which lead to $ 2g-1 $ parameters
with each parameter the dimension of the Poincare group $ ISO(2,1) $,
which is 6. Moreover, we need to quotient the phase space by the
conjugation of the Poincare group $ ISO(2,1) $, which lead to a
reduction of $ 6 $ parameters. Thus, the dimension of the phase space is:
\begin{align}
  \text{dim(Phase Space)} = (2g-1) \cdot 6 - 6 = 12g - 12
\end{align}

\rmk{
  Finally, we shall note that this investigation is not complete and
  rigorous and as Witten have noted that we need a field theoretical
  proof of the phase space. Moreover, though we can construct the solutions, we
  also have left the problem that the quotienting may lead to manifold with
  undesirable properties, though being flat.

  In the final section \cref{sec:Quantization of 3D Gravity} we will view the
  phase space from a more rigorous CS field theory approach.
}

\subsection{Action of Vierbein Formalism}

The Vielbein formalism of gravity can be written from a action principle.

\subsection{Review on Chern-Simons Theory}

\subsection{3D Gravity Action as Chern-Simons Theory}

\subsection{On-Shell Diffeomorphism as Gauge Transformation}
